%%%%%%%%%%%%%%%%%%%%%%%%%%%%%%%%%%%%%%%%%%%%%%%%%%%%%%%%%%%%
%%% ELIFE ARTICLE TEMPLATE
%%%%%%%%%%%%%%%%%%%%%%%%%%%%%%%%%%%%%%%%%%%%%%%%%%%%%%%%%%%%
%%% PREAMBLE 
\documentclass[9pt,lineno]{elife}
% Use the onehalfspacing option for 1.5 line spacing
% Use the doublespacing option for 2.0 line spacing
% Please note that these options may affect formatting.
% Additionally, the use of the \newcommand function should be limited.


\usepackage[version=4]{mhchem}
\usepackage{siunitx}
\usepackage{graphicx}
\DeclareSIUnit\Molar{M}

%%%%%%%%%%%%%%%%%%%%%%%%%%%%%%%%%%%%%%%%%%%%%%%%%%%%%%%%%%%%
%%% ARTICLE SETUP
%%%%%%%%%%%%%%%%%%%%%%%%%%%%%%%%%%%%%%%%%%%%%%%%%%%%%%%%%%%%
\title{This is the title}

\author[1*]{Joanna Jędrzejewska-Szmek}
\author[2*]{Kim T Blackwell}
\affil[1]{Institution 1}
\affil[2]{Institution 2}

\corr{email1@example.com}{JJS}
\corr{email2@example.com}{KTB}



%%%%%%%%%%%%%%%%%%%%%%%%%%%%%%%%%%%%%%%%%%%%%%%%%%%%%%%%%%%%
%%% ARTICLE START
%%%%%%%%%%%%%%%%%%%%%%%%%%%%%%%%%%%%%%%%%%%%%%%%%%%%%%%%%%%%

\begin{document}

\maketitle

\begin{abstract}
\end{abstract}


\section{Introduction (Level 1 heading)}




\section{Results (Level 1 heading)}



%% \begin{table}[bt]
%% \caption{\label{tab:example}Automobile Land Speed Records (GR 5-10).}
%% % Use "S" column identifier to align on decimal point 
%% \begin{tabular}{S l l l r}
%% \toprule
%% {Speed (mph)} & Driver          & Car                        & Engine    & Date     \\
%% \midrule
%% 407.447     & Craig Breedlove & Spirit of America          & GE J47    & 8/5/63   \\
%% 413.199     & Tom Green       & Wingfoot Express           & WE J46    & 10/2/64  \\
%% 434.22      & Art Arfons      & Green Monster              & GE J79    & 10/5/64  \\
%% 468.719     & Craig Breedlove & Spirit of America          & GE J79    & 10/13/64 \\
%% 526.277     & Craig Breedlove & Spirit of America          & GE J79    & 10/15/65 \\
%% 536.712     & Art Arfons      & Green Monster              & GE J79    & 10/27/65 \\
%% 555.127     & Craig Breedlove & Spirit of America, Sonic 1 & GE J79    & 11/2/65  \\
%% 576.553     & Art Arfons      & Green Monster              & GE J79    & 11/7/65  \\
%% 600.601     & Craig Breedlove & Spirit of America, Sonic 1 & GE J79    & 11/15/65 \\
%% 622.407     & Gary Gabelich   & Blue Flame                 & Rocket    & 10/23/70 \\
%% 633.468     & Richard Noble   & Thrust 2                   & RR RG 146 & 10/4/83  \\
%% 763.035     & Andy Green      & Thrust SSC                 & RR Spey   & 10/15/97\\
%% \bottomrule
%% \end{tabular}

%% \medskip 
%% Source: \url{https://www.sedl.org/afterschool/toolkits/science/pdf/ast_sci_data_tables_sample.pdf}

%% \tabledata{This is a description of a data source.}\label{tabdata:first}
%% \tablesrccode{This is a description of a source code.}\label{tabsrccode:first}

%% \end{table}

\subsection{Level 2 Heading}



\subsubsection{Level 3 Heading}



\paragraph{Level 4 Heading}


\section{Discussion}



\section{Methods and Materials}

Ryanodine receptor is a tertamer comprising four identical subunits, which forms a functional ion channel. Each ryanodine type 2 receptor (RyR2) subunit can bind calmodulin at three sites: CaMDB1, CaMDB2, CaMDB3 \citep{3_sites_ryr3cam}. Among these sites CaMBD1 has the lowest affinity for calmodulin \cite{3_sites_ryr3cam,ryrcam_sites_affinities} and its role is unknown, CaMBD2 and CaMBD3 have bindinging affinities in the range  of 1 uM to apoCaM to 50 nM to Ca-CaM \citep{3_sites_ryr3cam, low_threshold_soicrCaMBD1,xu_meissner_2004}. However, only calmodulin binding to CaMBD2 is resposible for the inhibitory effect of calmodulin on RyR2 \citep{YAMAGUCHI200323480,low_threshold_soicrCaMBD1} and we costructed a multistate model of RyR2-CaM complex with calmodulin bound at CaMBD2. CaM dissociates from RYR2 on the scale of minutes \citep{BALSHAW200120144}, which is increased to minutes for calcium-bound calmodulin \citep{BALSHAW200120144, ryrcam_dissociation_rate}.

Kinetic constant for calcium-activation of RyR2 Calcium activation of RyR2 unbound to calmodulin was fitted to RyR2 opening time course from \citep{ca_stoichiometry_ryr}, and verfied by evaluating the fit to open probability, mean open and closed time provided by \citep{xu_meissner_2004}. Calcium activation of CaM-bound RyR2 open probability was fitted to appropriate data from \citep{xu_meissner_2004} and the mean open and closed time was validated to best match data from \citep{xu_meissner_2004}.

We used the following activation scheme for RyR2 activation, based on the gating model of the Hodgik-Huxley potassium channel \citep{johnston_wu} and the kinetic model of RyR2 gating \citep{Zahradnikova_kon,ca_stoichiometry_ryr}
\begin{figure}
\includegraphics[width=0.95\linewidth]{RyR_Ca_activation_scheme.png}
\caption{\label{fig:RyR2activation}}
\end{figure}




%% \section{Some \LaTeX{} Examples}
%% \label{sec:examples}

%% Use \verb|\section| and \verb|\subsection| commands to organize your document. \LaTeX{} handles all the formatting automatically. Use \verb|\label| and \verb|\nameref| commands for cross-referencing sectional headings: the usual \verb|\ref| will not work, as this template uses unnumbered sectional headings.

%% \subsection{Figures and Tables}

%% Use the table and tabular commands for basic tables --- see \TABLE{example}, for example. 

%% You can upload a figure (JPEG, PNG or PDF) using the project menu. To include it in your document, use the \verb|\includegraphics| command as in the code for \FIG{view}. 

%% For a half-width figure or table with text wrapping around it, use 

%% \begin{verbatim}
%% \begin{wrapfigure}{l}{.46\textwidth}
%%   \includegraphics[width=\hsize]{...}
%%   \caption{...}\label{...}
%% \end{wrapfigure}
%% \end{verbatim}
%% %
%% as in \FIG{halfwidth}. For tables:

%% \begin{verbatim}
%% \begin{wraptable}{l}{.46\textwidth}{
%%   \begin{tabular}{...}
%%   ...
%%   \end{tabular}}
%%   \caption{...}\label{...}
%% \end{wraptable}
%% \end{verbatim}

%% Be careful with these, though, as they may behave strangely near page boundaries, sectional headings, or in the neighbourhood of lists or too many floats.

%% Labels for main videos can be added with \verb|\video| e.g.

%% \video{Ths is a description of a main video.\label{video:mv1}
%%   \videodata{This is a description of a video data source.}\label{viddata:first}
%%   \videodata{This is a description of another video data source.}\label{viddata:second}
%%   \videosrccode{This is a description of a video source code.}\label{vidsrccode:first}
%% }


%% Labels for video supplements can be added within \texttt{figure} environments, after the \texttt{caption}, using the \verb|\videosupp| command: see \VIDEOSUPP[view]{sv1} for an example.

%% If you use the following prefixes for your \verb|\label|:
%% %
%% \begin{description}
%% \item[Figures] \texttt{fig:}, e.g.~\verb|\label{fig:view}|
%% \item[Figure Supplements] \texttt{figsupp:}, e.g.~\verb|\label{figsupp:sf1}|\\
%% (we'll assume \texttt{figsupp:sf1} is a figure supplement of \texttt{fig:view} in our example)
%% \item[Figure source data] \texttt{figdata:}, e.g.~\verb|\label{figdata:first}|
%% \item[Figure source code] \texttt{figsrccode:},
%% \item[Videos] \texttt{video:}, e.g.~\verb|\label{video:mv1}|
%% \item[Video source data] \texttt{viddata:}, e.g.~\verb|\label{figdata:first}|
%% \item[Video source code] \texttt{vidsrccode:},
%% \item[Video supplements] \texttt{videosupp:}, e.g.~\verb|\label{videosupp:sv1}|
%% \item[Tables] \texttt{tab:}, e.g.~\verb|\label{tab:example}|
%% \item[Table source data] \texttt{tabdata:}, e.g.~\verb|\label{tabdata:first}|
%% \item[Table source code] \texttt{tabsrccode:},e.g.~\verb|\label{tabsrccode:first}|
%% \item[Equations] \texttt{eq:}, e.g.~\verb|\label{eq:CLT}|
%% \item[Boxes] \texttt{box:}, e.g.~\verb|\label{box:simple}|
%% \end{description}
%% %
%% you can then use the convenience commands \verb|\FIG{view}|, \verb|\FIGSUPP[view]{sf1}|, \verb|\TABLE{example}|, \verb|\EQ{CLT}|, \verb|\BOX{simple}|, \verb|\FIGDATA[view]{first}|, \verb|\FIGSRCCODE[view]{first}|, 
%% \verb|\TABLEDATA[example]{first}|, \verb|\TABLESRCCODE[example]{first}|,
%% \verb|\VIDEO{mv1}|,
%% \verb|\VIDEODATA[mv1]{second}|, 
%% \verb|\VIDEOSRCCODE[mv1]{first}|,
%% and also \verb|{\VIDEOSUPP}[view]{sv1}| \emph{without} the label prefixes, to generate cross-references 
%% \FIG{view}, \FIGSUPP[view]{sf1},  
%% \TABLE{example}, \EQ{CLT}, 
%% \BOX{simple}, 
%% \FIGDATA[view]{first}, 
%% \FIGSRCCODE[view]{first}, 
%% \TABLEDATA[example]{first}, 
%% \TABLESRCCODE[example]{first},
%% \VIDEO{mv1}, 
%% \VIDEODATA[mv1]{second}, 
%% \VIDEOSRCCODE[mv1]{first},
%% and \VIDEOSUPP[view]{sv1}. 
%% Alternatively, use \verb|\autoref| with the full label, e.g.~\autoref{first:app} (although this may not work correctly for figures and tables in the appendices or boxes nor supplements at present).

%% \begin{wrapfigure}{l}{.45\textwidth}
%% \includegraphics[width=\hsize]{frog}
%% \caption{A half-columnwidth image using wrapfigure, to be used sparingly. Note that using a wrapfigure before a sectional heading, near other floats or page boundaries is not recommended, as it may cause interesting layout issues. Use the optional argument to wrapfigure to control how many lines of text should be set half-width alongside it.}
%% \label{fig:halfwidth}
%% \end{wrapfigure}

%% Some filler text to sit alongside the half-width figure. \lipsum[1] \lipsum[2]

%% Really wide figures or tables, that take up the entire page, including the gutter space: use \verb|\begin{fullwidth}...\end{fullwidth}| as in \FIG{fullwidth}. And sometimes you may want to use feature boxes like \BOX{simple}.

%% \begin{figure}
%% \begin{fullwidth}
%% \includegraphics[width=0.95\linewidth]{elife-18156-fig2}
%% \caption{A very wide figure that takes up the entire page, including the gutter space. A very wide figure that takes up the entire page, including the gutter space. A very wide figure that takes up the entire page, including the gutter space. A very wide figure that takes up the entire page, including the gutter space. A very wide figure that takes up the entire page, including the gutter space. A very wide figure that takes up the entire page, including the gutter space.}
%% \label{fig:fullwidth}
%% \figsupp{There is no limit on the number of Figure Supplements for any one primary figure. Each figure supplement should be clearly labelled, Figure 1--Figure Supplement 1, Figure 1--Figure Supplement 2, Figure 2--Figure Supplement 1 and so on, and have a short title (and optional legend). Figure Supplements should be referred to in the legend of the associated primary figure, and should also be listed at the end of the article text file.}{\includegraphics[width=5cm]{frog}}
%% \end{fullwidth}
%% \end{figure}

%% \subsection{Citations}

%% LaTeX formats citations and references automatically using the bibliography records in your .bib file, which you can edit via the project menu. Use the \verb|\cite| command for an inline citation, like \cite{Aivazian917}, and the \verb|\citep| command for a citation in parentheses \citep{Aivazian917}. The LaTeX template uses a slightly-modified Vancouver bibliography style. If your manuscript is accepted, the eLife production team will re-format the references into the final published form. \emph{It is not necessary to attempt to format the reference list yourself to mirror the final published form.} Please also remember to \textbf{delete the line} \verb|\nocite{*}| in the template just before \verb|\bibliography{...}|; otherwise \emph{all} entries from your .bib file will be listed! 

%% \begin{featurebox}
%% \caption{This is an example feature box}
%% \label{box:simple}
%% This is a feature box. It floats!
%% \medskip

%% \includegraphics[width=5cm]{example-image}
%% \featurefig{`Figure' and `table' captions in feature boxes should be entered with \texttt{\textbackslash featurefig} and \texttt{\textbackslash featuretable}. They're not really floats.}

%% \lipsum[1]
%% \end{featurebox}

%% \subsection{Mathematics}

%% \LaTeX{} is great at typesetting mathematics $abc$. Let $X_1, X_2, \ldots, X_n$ be a sequence of independent and identically distributed random variables with $\text{E}[X_i] = \mu$ and $\text{Var}[X_i] = \sigma^2 < \infty$, and let
%% \begin{equation}
%% \label{eq:CLT}
%% S_n = \frac{X_1 + X_2 + \cdots + X_n}{n}
%%       = \frac{1}{n}\sum_{i}^{n} X_i
%% \end{equation}
%% denote their mean. Then as $n$ approaches infinity, the random variables $\sqrt{n}(S_n - \mu)$ converge in distribution to a normal $\mathcal{N}(0, \sigma^2)$.

%% \lipsum[3] 

%% \begin{figure}
%% \includegraphics[width=\linewidth]{elife-13214-fig7}
%% \caption{A text-width example.}
%% \label{fig:view}
%% %% If the optional argument in the square brackets is "none", then the caption *will not appear in the main figure at all* and only the full caption will appear under the supplementary figure at the end of the manuscript.
%% % 
%% \figsupp[Shorter caption for main text.]
%% {This is a supplementary figure's full caption, which will be used at the end of the manuscript. 
%%   \figsuppdata{A data source; see \url{https://doi.org/xxx}}
%%   \figsuppdata{Another data source.}
%%   \figsuppsrccode{And the source code.}}
%% {\includegraphics[width=6cm]{frog}}\label{figsupp:sf1}
%% % 
%% % 
%% \figsupp{This is another supplementary figure.}
%% {\includegraphics[width=6cm]{frog}}
%% % 
%% % 
%% \videosupp{This is a description of a video supplement.}\label{videosupp:sv1}
%% \figdata{This is a description of a data source.}\label{figdata:first}
%% \figdata{This is another description of a data source.}\label{figdata:second}
%% \figsrccode{This is a description of a source code.}\label{figsrccode:first}
%% \end{figure}

%% \subsection{Other Chemistry Niceties}

%% You can use commands from the \texttt{mhchem} and \texttt{siunitx} packages. For example: \ce{C32H64NO7S}; \SI{5}{\micro\metre}; \SI{30}{\degreeCelsius}; \SI{5e-17}{\Molar}

%% \subsection{Lists}

%% You can make lists with automatic numbering \dots

%% \begin{enumerate}
%% \item Like this,
%% \item and like this.
%% \end{enumerate}
%% \dots or bullet points \dots
%% \begin{itemize} 
%% \item Like this,
%% \item and like this.
%% \end{itemize}
%% \dots or with words and descriptions \dots
%% \begin{description}
%% \item[Word] Definition
%% \item[Concept] Explanation
%% \item[Idea] Text
%% \end{description}

%% Some filler text, because empty templates look really poorly. \lipsum[1]


\section{Acknowledgments}

Additional information can be given in the template, such as to not include funder information in the acknowledgments section.


\bibliography{elife-sample}

%%%%%%%%%%%%%%%%%%%%%%%%%%%%%%%%%%%%%%%%%%%%%%%%%%%%%%%%%%%%
%%% APPENDICES
%%%%%%%%%%%%%%%%%%%%%%%%%%%%%%%%%%%%%%%%%%%%%%%%%%%%%%%%%%%%

%% \appendix
%% \begin{appendixbox}
%% \label{first:app}
%% \section{Firstly}
%% \lipsum[1]

%% %% Sadly, we can't use floats in the appendix boxes. So they don't "float", but use \captionof{figure}{...} and \captionof{table}{...} to get them properly caption.
%% \begin{center}
%% \includegraphics[width=\linewidth,height=7cm]{frog}
%% \captionof{figure}{This is a figure in the appendix}
%% \end{center}

%% \section{Secondly}

%% \lipsum[5-8]

%% \begin{center}
%% \includegraphics[width=\linewidth,height=7cm]{frog}
%% \captionof{figure}{This is a figure in the appendix}
%% \end{center}

%% \end{appendixbox}

%% \begin{appendixbox}
%% \includegraphics[width=\linewidth,height=7cm]{frog}
%% \captionof{figure}{This is a figure in the appendix}
%% \end{appendixbox}
\end{document}
